\section{Repository records cited in this note}
\label{app:provenance}

Several statements in this note rest on records in the analysis repository rather than on anything
a reader can recompute from the published inputs. Those records are collected here rather than
inlined, for two reasons. A repository path is not usable by a reader outside the collaboration, so
in the body it interrupts a physics argument with a string that cannot be followed; and gathering
them in one table means a path that moves has one place to be corrected instead of five.
All paths below are relative to the code repository, \repolink, and are set breakably so they stay
inside the text block.

\textbf{This appendix is not the reproduction path.} What an external reader would need in order to
recompute the quoted numbers --- the arrays, edges, orderings, masks, units, projection matrices and
generator predictions, with one worked example --- is specified separately in
App.~\ref{app:release}, together with a measured statement of which of those objects exist today and
which do not.

\textbf{This appendix is provenance, not evidence.} Nothing here is a result, and no statement in
the body becomes stronger because its record is listed. Where a claim depends on a record whose own
limitations matter---most importantly the adopted covariance's four measurements
(\S\ref{sec:eavailw-adopted})---those limitations are stated at the claim, not delegated here.

The tags are stable; the paths are current as of 2026-10-04. Rows are ordered by the
section that cites them. Table~\ref{tab:prov-paths} lists records with a repository path;
Table~\ref{tab:prov-trackers} lists tracker rows and code identifiers that have no
document of their own.

{\small
\begin{longtable}{@{}llp{0.50\textwidth}@{}}
\caption{Repository records with a path.}\label{tab:prov-paths}\\
\toprule
Tag & Cited at & Record \\
\midrule
\endfirsthead
\toprule
Tag & Cited at & Record \\
\midrule
\endhead
\midrule \multicolumn{3}{r}{\emph{continued\ldots}}\\ \endfoot
\bottomrule
\endlastfoot
\texttt{coverage-toys} & \S\ref{sec:coverage} &
  see below \\
\texttt{flux-repair} & \shortstack[l]{\S\ref{sec:3d-models}, \S\ref{sec:eavailw-band},\\
  App.~\ref{app:release-schema}, App.~\ref{app:history}} &
  see below \\
\texttt{q3-excess} & \S\ref{sec:4d} &
  \texttt{nd-unfolding/}\allowbreak\texttt{q3\_excess\_projection.py} (comparator: Tune~v1 prior,
  weights \texttt{w\_truth}); produces Fig.~\ref{fig:q3excess} \\
\texttt{eavailw-excess} & Fig.~\ref{fig:eavailWexcess} &
  \texttt{nd-unfolding/}\allowbreak\texttt{excess\_eavail\_W.py} (cell-integrated differences;
  comparator weights \texttt{w\_truth}) \\
\texttt{z-adoption} & \S\ref{sec:eavailw-construction}, App.~\ref{app:release-audit} &
  see below \\
\texttt{retracted-values} & \S\ref{sec:pet}, App.~\ref{app:history} &
  \repopath{docs/orchestration/INDEX-retracted-and-superseded-values.md} \\
\texttt{pet-pairing} & \S\ref{sec:pet}, \S\ref{app:cstatlimit} &
  see below \\
\texttt{pet-closure-legacy} & \S\ref{sec:pet-abs} &
  ruling of 2026-08-21;
  \repopath{docs/orchestration/DETERMINATION-20260819-lanec-petclosure-is-legacy.md} \\
\texttt{pet-projection-figure} & \S\ref{sec:pet-projection} &
  see below \\
\texttt{variance-sources} & Table~\ref{tab:variance-sources} (\S\ref{sec:variance}) &
  see below \\
\texttt{rank-inversion} & \S\ref{sec:nd-declarations} &
  \repopath{docs/orchestration/RANK-AND-INVERSION-20260810.md}; the predeclared rank cut itself is
  a statistics decision recorded in \texttt{docs/}\allowbreak\texttt{STATVAL\_REPAIR.md} \\
\texttt{closure-scripts} & \S\ref{sec:closure} &
  \texttt{uq/closure/}\allowbreak\texttt{closure\_truth\_reweight.py} (\S\ref{sec:closure-tr});
  \texttt{uq/closure/}\allowbreak\texttt{closure\_hidden\_var.py} (\S\ref{sec:closure-hv});
  \texttt{uq/closure/}\allowbreak\texttt{closure\_alt\_model.py} (\S\ref{sec:closure-alt}) \\
\texttt{cstat-recompute} & \S\ref{app:cstatlimit} &
  see below \\
\texttt{negweight-toy} & \S\ref{app:negweight} &
  \repopath{docs/orchestration/receipts/RECEIPT-negweight-toy-20260821.json} \\
\texttt{negweight-hpss} & \S\ref{app:negweight} &
  see below \\
\texttt{negweight-audit} & \S\ref{app:negweight} &
  \repopath{docs/orchestration/AUDIT-FINDINGS-20260731.md}, tier~A, finding J22 \\
\texttt{gbdt-closeout} & \S\ref{app:negweight} &
  \repopath{docs/orchestration/RUNBOOK-20260807-gbdt-closeout.md} \S2.1--2.2 \\
\texttt{oi6-oi8-oi126} & \S\ref{app:negweight} &
  see below \\
\texttt{p4-verifier} & \S\ref{app:negweight} &
  \repopath{docs/orchestration/runs/standard-p4-verifier/} \\
\texttt{fps-battery} & \S\ref{app:negweight} &
  see below \\
\texttt{z-projection} & App.~\ref{app:release-identities} &
  see below \\
\texttt{z-assembly} & App.~\ref{app:release-audit} &
  \texttt{nd-unfolding/}\allowbreak\texttt{z\_assembly.py} (implements
  Eqs.~\eqref{eq:zassembly}--\eqref{eq:zinflation}); the same algebra is asserted independently
  of the producer \\
\texttt{z-conditioning} & App.~\ref{app:release-audit} &
  see below \\
\end{longtable}
}

\noindent Tags marked ``see below'' carry records or context too long for the
column above, so they are given here at full width.

\begin{description}[leftmargin=0pt,itemsep=3pt]
  \item[\texttt{coverage-toys}.] Launcher \texttt{sbatch\_coverage\_toys\_MEFHC\_200.sh}; rollup
    \texttt{uq/coverage\_toys.py}; summary \texttt{uq/coverage/}\allowbreak\texttt{coverage\_summary.txt}.
    Each toy runs the driver with \texttt{-{}-closure} and \texttt{-{}-bootstrap-seed $k$}; the
    stored truth histogram is \texttt{hTruthXSec2D}, fluctuated through \texttt{w\_truth}.
  \item[\texttt{flux-repair}.] \texttt{KNOWN\_ISSUES.md} rows 82 (GiBUU sample stops at
    $E_\nu\approx\SI{20}{GeV}$) and 83 (GENIE/NuWro flux mis-sampling); repair and rebuild of
    2026-09-27, \texttt{VALIDATION\_LEDGER.md} rows \texttt{VL156}--\texttt{VL160} (\texttt{VL159}:
    3D integrated-rate shares; \texttt{VL160}: Valencia-2p2h fill), superseding
    \texttt{VL35}--\texttt{VL38}; receipt
    \repopath{docs/orchestration/state/s5p/stage7/generator-context/generator-context-receipt.json}.
  \item[\texttt{z-adoption}.] The adoption itself at
    \repopath{docs/orchestration/DECISION-20260920-joseph-adopts-z-cv-under-the-6.4-exception.md};
    the seed measurements at
    \repopath{docs/orchestration/EVIDENCE-20260920-sproj-resolution-floor-and-seed-effect.md}
    and \repopath{docs/orchestration/state/SEED-EFFECT-20260920.json}; and the
    withdrawal of the lower-bound reading at
    \repopath{docs/orchestration/CORRECTION-20260920-lower-bound-inference-withdrawn.md}.
    Adoption recorded 2026-09-20; the predeclared magnitude criterion is \emph{cause~3}'s
    $M(\mathrm{ii})$ in the internal dependency register, computed by the 2026-09-20 campaign
    for the two companion products only.
  \item[\texttt{pet-pairing}.] \texttt{OI-126} in \texttt{docs/OPEN\_ITEMS.md}.
    The 2026-08-20 ruling: the central/statistical pairing is declined, no PET
    total covariance is adopted for publication, PET remains a
    diagnostic and method-development result, and reconsideration requires
    estimator-equivalence \emph{and} coverage evidence --- coverage being a
    different object from verifying the construction.
    Ratification rests with the estimator's owner and the covariance-construction reviewer, and
    the open item recording the question (\texttt{OI-126}) is not closed by
    App.~\ref{app:cstatlimit}.
  \item[\texttt{pet-projection-figure}.] \texttt{ISSUE-59} in
    \texttt{KNOWN\_ISSUES.md} (fixed 2026-09-23; \texttt{pet\_cloud\_projection\_xsec} regenerated) and the corrected values at
    \texttt{nd-unfolding/products/pet/fullcloud/pointcloud\_projection\_summary.json},
    which holds the three $\Eavail$ projections (\texttt{pet\_cloud\_hascloud},
    \texttt{pet\_stored\_hascloud}, \texttt{pet\_stored\_full}); method and
    reproduction at \texttt{nd-unfolding/pet/POINTCLOUD\_PROJECTION.md}.
  \item[\texttt{variance-sources}.] Code pointers, one per row of
    Table~\ref{tab:variance-sources}. Training-seed variation:
    \texttt{seedscan/analyze\_seedscan.py}; \texttt{uq/seedscan\_lgbm\_ml/
    uq\_covariance\_ml.root}. Poisson bootstrap: \texttt{unfold\_2d\_omnifold\_unbinned.py}
    (\texttt{-{}-bootstrap-seed}, \texttt{-{}-seed}); \texttt{uq/analyze\_uq.py};
    \texttt{uq/bootstrap\_MEFHC\_300/ uq\_covariance\_boot300.root}. Universes:
    \texttt{runEventLoopOmniFold.cpp} (\texttt{MNV101\_DUMP\_UNIVERSES});
    \texttt{unfold\_2d\_omnifold\_unbinned.py -{}-universe BAND:IDX};
    \texttt{uq/analyze\_universes.py} (\texttt{-{}-add-norm 0.014}). MAT reference:
    \texttt{uq/verify\_matcorr\_vs\_mnvh1d.py}. Block sum:
    \texttt{compare\_to\_paper\_fullcov.py} (\texttt{-{}-omnifold-cov}).
  \item[\texttt{cstat-recompute}.] Every figure in App.~\ref{app:cstatlimit} is derived from
    committed artifacts, and the per-cell arrays behind the aggregates are published with them:
    the family spread and its per-replica values in
    \texttt{state/probe-oi126-band-Rpush-sigma-20260815}; the training-versus-extraction split,
    its control region, and the thresholds fixed before that comparison ran in
    \texttt{state/probe-oi126-push-vs-extraction-20260815} and the predeclaration beside it; the
    $\mathrm{Poisson}(1)$ verification of the targets in
    \texttt{state/RECEIPT-20260815-oi126-mechanism-narrowing} and
    \texttt{state/probe-oi126-zero-fraction-per-column-20260815}; the band geometry, acceptance,
    share of integral and the external comparison in
    \texttt{state/RECEIPT-20260815-cstat-tail-geometry-and-weighting-correction}; and the two
    per-cell relative standard deviations tabulated there as fields of
    \texttt{state/p5a-nominal-vs-cstat-family-percell-20260815}, recomputable from that receipt's
    own shipped per-cell array. No share-of-total figure from that receipt is used
    (App.~\ref{app:history}). The members' input file is measured in
    \texttt{state/gate5-source-npz-verified-20260813}, which should be cited for that purpose in
    preference to the members' own recorded input digest.
  \item[\texttt{negweight-hpss}.]
    \repopath{docs/orchestration/RECEIPT-20260821-negweight-hpss-durability.md}.
    The real-data products behind the twelve
    production values are archived and were restored end to end from tape under a tested
    recovery route; the synthetic-toy values come from a committed, fully deterministic
    script that reproduces all four on re-run (\texttt{negweight-toy}). Two provenance
    limits travel with the production numbers: the driver was uncommitted when those
    outputs were written, so the executing revision is corroborated by a version-distinct
    log string rather than pinned by a digest, and no output carries a producing job
    identifier, which no later run can supply.
  \item[\texttt{oi6-oi8-oi126}.]
    \repopath{docs/orchestration/DECISION-20260815-joseph-oi6-oi8-oi126.md}. The purity
    footing of the standard chain was raised as a question --- whether the locked
    instruction that $N$-D production runs \texttt{negweight-refined} reaches the
    standard-phase-space chain, or only its extended-fiducial form --- and answered in
    favour of purity on 2026-08-07 (\texttt{gbdt-closeout}) and again, after the
    alternative was costed, on 2026-08-15.
  \item[\texttt{fps-battery}.] FPS ME~FHC battery of 2026-06-10, scheduler job
    \texttt{54244120} (\texttt{VALIDATION\_LEDGER.md}, ``FPS MEFHC battery''). The
    negative-weight driver code did not land until 2026-07-11 (commit \texttt{cf8a4a6})
    and the \texttt{-{}-bkg-mode negweight-refined} launcher flag not until 2026-07-20.
    Both FPS unfold launchers omitted \texttt{-{}-bkg-mode} entirely, so the ten endpoint
    unfolds ran on the driver default \texttt{-{}-bkg-mode=purity}; they are recorded as
    purity controls at
    \texttt{nd-unfolding/uq\_fps/corrected/}\allowbreak\texttt{FPS\_UQ\_CORRECTED\_STATE.md}.
  \item[\texttt{z-projection}.]
    \repopath{docs/orchestration/state/PROJ-20260920-m1-publication-receipt.json}
    and \repopath{docs/orchestration/state/PROJ-20260920-binding-check.json}.
  \item[\texttt{z-conditioning}.] Cutoff scan:
    \texttt{PLAN-20260918-scalar5d-publication-completion.md} \S16.1a (diagnostic projection, same
    $M$); independent scan of the publication bytes: \texttt{VALIDATION\_LEDGER.md} row
    \texttt{VL145} (publication and diagnostic $42\times42$ matrices bitwise identical). Default
    cutoff $10^{-12}$ is hardcoded in \texttt{project\_cov\_nd.py}.
\end{description}

{\small
\begin{longtable}{@{}llp{0.25\textwidth}p{0.30\textwidth}@{}}
\caption{Tracker rows and code identifiers: \texttt{KNOWN\_ISSUES.md} rows,
\texttt{VALIDATION\_LEDGER.md} rows (VL), open items (OI), issues, audit findings,
gates and receipt fields, commits and job identifiers, and code symbols.}\label{tab:prov-trackers}\\
\toprule
Tag & Cited at & Record & Content \\
\midrule
\endfirsthead
\toprule
Tag & Cited at & Record & Content \\
\midrule
\endhead
\midrule \multicolumn{4}{r}{\emph{continued\ldots}}\\ \endfoot
\bottomrule
\endlastfoot
\texttt{pet-code} & \S\ref{sec:pet-inputs}, \S\ref{sec:pet-abs} &
  \texttt{nd-unfolding/pet/}\allowbreak\texttt{minerva\_pet\_dataloader.py}; \texttt{MultiFold.reweight}
  (\texttt{omnifold\_nn/omnifold/}\allowbreak\texttt{omnifold.py}); \texttt{extract\_cross\_section\_nd}
  (\texttt{nd-unfolding/xsec\_nd.py}) &
  the dataloader slices off PDG before classification; the full-sample reweight;
  the standard extraction \\
\texttt{cstat-family} & \S\ref{sec:pet-syst}, \S\ref{app:cstatlimit} &
  \texttt{VL132} (the Gate-5 family); \texttt{CSTAT-R7} &
  the full-event extended-fiducial 50-member replica family; the design document specifying
  $N=100$ labels a family of this size \texttt{INSUFFICIENT}, and because that document's
  supersession is unratified, the label stands in the record as written \\
\texttt{pet-gate6} & \S\ref{sec:pet-full-event-status} &
  Gate~6 (blocked); receipt prohibitions \texttt{do\_not\_select\_passing\_subset},
  \texttt{do\_not\_construct\_C\_ML}, \texttt{do\_not\_move\_central},
  \texttt{do\_not\_start\_leg\_2}, \texttt{do\_not\_retry\_unchanged} &
  construction of a PET ML-variance component is not authorized to start \\
\texttt{pot-ratio} & \S\ref{sec:tension} (footnote) &
  \texttt{KNOWN\_ISSUES} \texttt{J36} &
  one global $\mathrm{POT}_{\mathrm{data}}/\mathrm{POT}_{\mathrm{MC}}$ ratio rather than a
  per-playlist ratio \\
\texttt{bkg-mode-pins} & \S\ref{sec:negweight-footing} &
  \texttt{STANDARD\_BKG\_MODE} (\texttt{nd-unfolding/p4\_lib.py});
  \texttt{PUBLICATION\_BKG\_MODE} (\texttt{nd-unfolding/}\allowbreak\texttt{fps\_provenance.py}) &
  standard-chain purity footing; extended-fiducial designated footing \\
\end{longtable}
}

\paragraph{One record is deliberately \emph{not} moved here.} The adopted five-axis covariance is
identified in the body by its \texttt{sha256} (\S\ref{sec:eavailw-adopted}), and that digest stays
in the body. It is not a routing pointer: it is the identity of the object every quoted
uncertainty is computed from, it is checkable by anyone holding the file, and the pairing between
the covariance and the central values is established \emph{by} that digest rather than by numerical
agreement. Moving it to an appendix would relocate the one identifier a reader needs in order to
know which object the numbers describe.
